\hypertarget{congruences}{%
\section{Distribution in residue classes}\label{congruences}}

We prove the negative-parity estimate used in
\protect\hyperlink{sieve}{the sieve argument}. All weights, prime bands,
and normalizing factors are those of \protect\hyperlink{proof}{the main
proof}. Put \(x=e^L\), \(Y=\lfloor2x\rfloor\), and
\(X=X_{-1}=\sum_{u\le Y}w_{-1}(u)\). For odd squarefree \(d\), set

\[
r(d)=\sum_{\substack{u\le Y\\u\equiv-2\pmod d}}w_{-1}(u)
       -\frac X{\varphi(d)}.
\]

For fixed \(0<\vartheta<1/2\) and every \(D>0\), the result is

\[
\sum_{\substack{d\le x^\vartheta\\d\text{ odd, squarefree}}}|r(d)|
 \ll_{\Psi,\vartheta,D}xL^{-D}+XL^{-18}. \tag{G1}
\]

The proof below concerns \(w_{-1}\). No distribution assertion for the
other parity sign is needed here.

\Needspace{6\baselineskip}

\hypertarget{congruences-orthogonality-and-the-nonunit-correction}{%
\subsection{Orthogonality and the nonunit
correction}\label{congruences-orthogonality-and-the-nonunit-correction}}

Extend Dirichlet characters by zero on nonunits, and write

\[
S_d(\chi)=\sum_{u\le Y}w_{-1}(u)\chi(u),\qquad
N(d)=\sum_{\substack{u\le Y\\(u,d)>1}}w_{-1}(u).
\]

Since \(-2\) is a unit modulo an odd \(d\), character orthogonality
gives

\[
\varphi(d)r(d)
 =\sum_{\chi\ne\chi_0}\chi((-2)^{-1})S_d(\chi)-N(d).
 \tag{G2}
\]

The correction \(N(d)\) occurs because the principal character sums to
\(X-N(d)\), rather than \(X\). In particular,

\[
|r(d)|\le\frac1{\varphi(d)}
 \left(\sum_{\chi\ne\chi_0}|S_d(\chi)|+N(d)\right). \tag{G3}
\]

Let

\[
H(N)=\sum_{\substack{d\le N\\d\text{ squarefree}}}\frac1{\varphi(d)}.
\]

For a fixed \(u\), assign a noncoprime squarefree \(d\) to any prime
\(p\mid(u,d)\) and write \(d=pe\). Dropping the conditions \((p,e)=1\)
and \(e\le N/p\) gives

\[
\sum_{\substack{d\le N\\d\text{ squarefree}\\(u,d)>1}}
 \frac1{\varphi(d)}
 \le H(N)\sum_{p\mid u}\frac1{p-1}. \tag{G4}
\]

The squarefree totient sum satisfies
\(H(\lfloor x^\vartheta\rfloor)\le eL\) eventually. Every prime divisor
of a supported \(u\le2x\) is at least \(L^{20}\). Also
\(2^{\omega(u)}\le u\), where \(\omega\) counts distinct prime divisors.
Therefore, eventually,

\[
\sum_{p\mid u}\frac1{p-1}
 \le\frac{2\omega(u)}{L^{20}}
 \le\frac4{\log2}L^{-19}.
\]

Multiply (G4) by \(w_{-1}(u)\ge0\) and interchange the finite sums. This
proves

\[
\sum_{\substack{d\le x^\vartheta\\d\text{ odd, squarefree}}}
 \frac{N(d)}{\varphi(d)}
 \le\frac{4e}{\log2}XL^{-18}. \tag{G5}
\]

\Needspace{6\baselineskip}

\hypertarget{congruences-exposing-one-prime-variable}{%
\subsection{Exposing one prime
variable}\label{congruences-exposing-one-prime-variable}}

Reserve one labelled slot in \(Q_0\), whose prime values satisfy

\[
x^\alpha\le q\le x^{6\alpha/5},\qquad \alpha=1/44.
\]

The ordered-slot expansion of the ordinary coefficient gives an exact
convolution

\[
w_{-1}(u)=\Psi(u/x)
 \sum_{\substack{q\in Q_0\\q\mid u}}b_L(u/q). \tag{G6}
\]

Here \(b_L\) is the sum over the remaining ordinary slots and all marked
group factors, with the original factorial normalization and negative
parity projection. Deleting the slot changes the number of slot
variables, not that normalization. If a prime occurs several times, its
different slot assignments remain in this finite expansion. After
factorial normalization the ordinary-slot coefficient is at most one,
and the marked factors are bounded by a fixed power of \(L\). Hence, for
a fixed \(c\ge0\),

\[
0\le b_L(R)\le L^c. \tag{G7}
\]

If \(\chi\) modulo \(d\) is induced by a primitive character \(\chi^*\)
of conductor \(f\mid d\), put \(h=d/f\). Inserting (G6) gives the exact
sum

\[
S_d(\chi)=\sum_{\substack{q\in Q_0,\ R\ge1\\(q,h)=(R,h)=1}}
 b_L(R)\Psi(qR/x)\chi^*(q)\chi^*(R). \tag{G8}
\]

The two coprimality masks are retained. Since \(d\) is squarefree,
\((f,h)=1\) and \(\varphi(d)=\varphi(f)\varphi(h)\). Summing the
complementary level \(h\) costs at most
\(H(\lfloor x^\vartheta\rfloor)=O(L)\).

Partition the two variables into binary dyadic intervals \(q\asymp P\),
\(R\asymp U\). The support \(x\le qR\le2x\) leaves \(O(L)\) pairs of
scales, and they satisfy

\[
P\le x^{6\alpha/5},\qquad U\le4x^{1-\alpha},\qquad
x/4<PU\le2x. \tag{G9}
\]

It remains to bound the primitive-character sums on each such rectangle.

\Needspace{6\baselineskip}

\hypertarget{congruences-small-and-large-conductors}{%
\subsection{Small and large
conductors}\label{congruences-small-and-large-conductors}}

The two analytic inputs are the prime Siegel--Walfisz estimate and the
multiplicative large sieve. In the forms used here, Siegel--Walfisz
gives arbitrary logarithmic saving for prime partial sums twisted by a
nonprincipal character of modulus at most a fixed power of \(L\). The
large sieve bounds the primitive-character quadratic mean of a sequence
on an interval of length \(N\) by its squared \(\ell^2\) norm times
\(Q^2+\pi N\), with the usual \(f/\varphi(f)\) conductor weight for
\(f\le Q\).

Choose a conductor cutoff \(2^K\le L^{C_*}\). For \(f\le2^K\), apply
Siegel--Walfisz to the \(q\) variable in (G8). Abel summation
incorporates \(\Psi(qR/x)\); its supremum and total variation are fixed
constants. The excluded primes \(q\mid h\) number at most
\(\vartheta/\alpha\), since \(q\ge x^\alpha\) and \(h\le x^\vartheta\).
Thus the small-conductor contribution on one retained rectangle is at
most

\[
L^{C_*+c}\left(
 2xB_\Psi L^{-A}V_\Psi
 +4x^{1-\alpha}\frac\vartheta\alpha S_\Psi\right), \tag{G10}
\]

where \(B_\Psi,S_\Psi,V_\Psi\) are fixed nonnegative constants from the
prime partial-sum estimate and the cutoff. The exponent \(A\) may be
chosen arbitrarily large once \(C_*\) is fixed.

For \(F<f\le2F\), use Fourier inversion in the logarithmic variable to
separate \(\Psi(qR/x)\). If \(\psi(t)=\Psi(e^t)\), write
\(M_\Psi=\int_{\mathbb R}|\widehat\psi(s)|\,ds<\infty\). At each
frequency, Cauchy--Schwarz in the primitive characters, followed by the
large sieve for the two factors, yields

\[
\mathcal B_F\le M_\Psi L^c
 \sqrt{\frac{(4F^2+\pi P)P}{F}}
 \sqrt{\frac{(4F^2+\pi U)U}{F}}. \tag{G11}
\]

Here \(\mathcal B_F\) is the sum of the absolute bilinear character sums
weighted by \(1/\varphi(f)\), for the fixed complementary level \(h\)
and rectangle. The masks in (G8) only decrease the coefficient norms.
Expanding the right side with \(\sqrt{a+b}\le\sqrt a+\sqrt b\), using
(G9), and summing over \(2^K<F\le2^T\) gives

\[
\sum_F\mathcal B_F\le M_\Psi L^c\left(
8\cdot2^T x^{1/2}
+12\sqrt\pi\,T x^{1-\alpha/2}
+\frac{4\pi x}{2^K}\right). \tag{G12}
\]

For a requested exponent \(D\), take

\[
k=c+D+2,\qquad C_*=k+1,\qquad A=C_*+c+D+2,
\]

and choose integer \(K,T\) with

\[
L^k\le2^K\le L^{k+1},\qquad
x^\vartheta\le2^T\le2x^\vartheta.
\]

The two remaining costs, the complementary levels and the rectangle
count, are each \(O(L)\). They turn the first term of (G10) and the last
term of (G12) into \(O(xL^{-D})\). All other terms contain a fixed
negative power of \(x\) after division by \(x\): the first term in (G12)
uses exactly \(\vartheta<1/2\). These powers absorb every fixed
logarithmic factor. We conclude that

\[
\sum_{\substack{d\le x^\vartheta\\d\text{ odd, squarefree}}}
 \frac1{\varphi(d)}\sum_{\chi\ne\chi_0}|S_d(\chi)|
 \ll xL^{-D}. \tag{G13}
\]

Equations (G3), (G5), and (G13) prove (G1). At \(D=230\), the lower mass
bound \(X\gg xL^{-221}\) makes its right side \(o(X/L)\).
